\input{Iteration-V2/Chapter3/chapter-03-23}

\subsection{Revision 24 requirements and trust-store threat model}

Revision 24 adds eight requirements. \textbf{TS1} requires a fixed canonical schema whose anchor and revocation ordering cannot alter the signed payload. \textbf{TS2} requires one structurally valid Ed25519 offline root that is valid both at envelope issuance and assessment. \textbf{TS3} forbids an offline-root identifier from also identifying a reviewer anchor and requires different issuer and approver identities. \textbf{TS4} requires explicit, dated revocations within the envelope validity window. \textbf{TS5} requires sequence one to have no predecessor and later sequences to advance exactly one step while binding the accepted predecessor digest. \textbf{TS6} rejects rollback, same-sequence forks, gaps, missing history, invalid signatures, freeze through expiry, and future activation. \textbf{TS7} persists the highest accepted sequence and digest before exposing reviewer anchors and removes all effective anchors if persistence fails. \textbf{TS8} requires the interface to expose both the production provisioning state and local-history limitations.

The attacker model includes mutation of signed fields, substitution of an unknown or malformed root, replay of an older valid envelope, same-sequence replacement, omission of intermediate envelopes, predecessor substitution, and continued use after expiry. The design does not resist an attacker who can replace application code or roots, compromise the device, erase application storage, control the release channel, or subvert the offline ceremony. Those are residual distribution and platform threats, not cases silently counted as passing.


\subsection{Revision 25 witness requirements and attacker model}

Revision 25 adds seven requirements. \textbf{WR1} requires a fixed receipt schema and canonical field order. \textbf{WR2} rejects thresholds below two or above the configured anchor count. \textbf{WR3} requires unique witness owners, key identifiers, and public keys, with identifiers and public-key material distinct from roots and reviewers. \textbf{WR4} requires valid Ed25519 material, coherent validity and revocation dates, and validity at both witnessing and assessment. \textbf{WR5} binds every receipt to the exact envelope digest, sequence, timely witness date, and expected application release identity. \textbf{WR6} rejects the entire supplied set on malformed, duplicate, unknown, revoked, future-dated, wrongly bound, or incorrectly signed receipts. \textbf{WR7} exposes no reviewer anchors and advances no rollback state until quorum passes.

The extended attacker model includes a compromised root acting alone, receipt replay against another envelope or release, duplicate counting, threshold reduction to one, witness/root identity collision, signature mutation, future dating, and use of a revoked witness. It does not claim resistance to colluding or fictitious witnesses, replacement of application code or configured anchors, release-channel compromise, public split views, or device compromise.


\subsection{Revision 26 human-decision requirements}

Five additional requirements govern the decision pause. \textbf{DR1} leaves all evidence visible without requiring an answer. \textbf{DR2} distinguishes a technical signal from both a GDPR violation and developer intent. \textbf{DR3} never treats selected checkboxes as proof of comprehension or legal correctness. \textbf{DR4} retains answers only in component memory and clears them when the card closes. \textbf{DR5} exposes radio, checkbox, selected-state, 44-pixel target, and polite live-feedback semantics to accessibility services.

The corresponding misuse model includes automation bias, premature consequential action, coerced agreement, dark-pattern gating, telemetry disguised as a comprehension exercise, and synthetic-study results presented as participant evidence. Revision 26 reduces the first two through corrective friction and blocks the latter four by preserving evidence access, collecting nothing in the App, and attaching an explicit evidence-class field to research input. It cannot guarantee attention, comprehension, emotional safety, or proportionate behaviour.
